\par\noindent\textbf{(1)} 在δ点x0波函数连续，导数跳跃ψ'(x0+)−ψ'(x0−)=−(2mg/ℏ²)ψ(x0)。势关于原点对称，束缚态可取偶、奇宇称；在|x|>a用指数衰减e\textasciicircum{}\{−κ|x|\}，在0<x<a分别取coshκx或sinhκx。
\par\noindent\textbf{(2)} 将内区偶态coshκx与外区指数衰减匹配，得到κ=(mg/ℏ²)(1+e\textasciicircum{}\{−2κa\})；奇态sinhκx则给κ=(mg/ℏ²)(1−e\textasciicircum{}\{−2κa\})。偶态总有唯一正κ；奇态可能不存在，因为其方程在κ→0附近需要足够大的斜率。
\par\noindent\textbf{(3)} 记κ0=mg/ℏ²，若κ0a≫1，两根κ\_±≈κ0[1±e\textasciicircum{}\{−2κ0a\}]，对应能级E\_±≈−ℏ²κ0²/(2m)[1±2e\textasciicircum{}\{−2κ0a\}]。偶态更负、奇态稍高，能量间距ΔE≈(2ℏ²κ0²/m)e\textasciicircum{}\{−2κ0a\}，来自两局域井的隧穿重叠。
\par\noindent\textbf{(4)} 奇态方程κ=κ0(1−e\textasciicircum{}\{−2κa\})在κ→0的斜率条件为2κ0a>1，即2mga/ℏ²>1；等号时奇态恰碰连续谱阈值而非可归一束缚态。a→0时两个δ合成强度2g的单δ，偶态κ→2κ0、奇态消失。
